% =========================================================================
% Direction Error Bound (DEB) — Standalone formulation for review
% To be integrated into IEEE_TCOM_RV1.tex after verification
% =========================================================================

\documentclass[journal]{IEEEtran}
\usepackage{amsmath,amsfonts,amssymb}
\usepackage{bm}

\begin{document}

\section*{Direction Error Bound (DEB) — Formulation}

\subsection*{Objective}
Derive the Cram\'{e}r--Rao lower bound for the direction vector 
$\mathbf{n}_d \in \mathbb{S}^2$ estimated from the $K$ direction-finding 
measurements only (excluding the distance-recovery measurement $K{+}1$).
This bound, termed the Direction Error Bound (DEB), quantifies the minimum 
achievable RMSE of $\|\hat{\mathbf{n}}_d - \mathbf{n}_d\|$ for any unbiased 
direction estimator.

% -----------------------------------------------------------------
\subsection*{1. Direction-Finding Observation Model}

During the direction-finding stage, the sample-mean optical power at the 
$i$-th transmitter orientation is:
\begin{equation}
  \bar{P}_{r,i} \sim \mathcal{N}\!\bigl(\mu_i,\;\sigma^2/N\bigr),
  \quad i = 1,\dots,K,
  \label{eq:deb_obs}
\end{equation}
where the noise-free mean power is:
\begin{equation}
  \mu_i = \frac{C}{d^2}\,\cos^m\!\phi_i\;\cos\psi
        = \eta\,(\mathbf{n}_{t,i}\cdot\mathbf{n}_d)^m,
  \label{eq:deb_mu}
\end{equation}
with $C = P_t(m{+}1)A_{\mathrm{det}}/(2\pi)$ and the \emph{nuisance amplitude}
\begin{equation}
  \eta \;\triangleq\; \frac{C\cos\psi}{d^2},
  \label{eq:deb_eta}
\end{equation}
which absorbs the unknown distance $d$ and the receiver-dependent factor 
$\cos\psi = \mathbf{n}_r\!\cdot\!(-\mathbf{n}_d)$.
The direction $\mathbf{n}_d$ and the scalar $\eta$ are the only unknowns; 
all other quantities ($\mathbf{n}_{t,i}$, $m$, $\sigma^2$, $N$) are known.

\emph{Remark.}  
Because $\eta$ absorbs $\cos\psi$, the direction-finding model 
\eqref{eq:deb_mu} does \emph{not} depend on $\mathbf{n}_r$.  
This is consistent with Proposition~1: the power ratios 
$\beta_i = (\mu_i/\mu_1)^{1/m}$ that underlie GLS/WLS are 
$\mathbf{n}_r$-independent.

% -----------------------------------------------------------------
\subsection*{2. Spherical Parameterization}

The unit direction vector $\mathbf{n}_d$ lies on $\mathbb{S}^2$ and is 
parameterized by two angles $(\theta_d, \varphi_d)$:
\begin{equation}
  \mathbf{n}_d(\theta_d,\varphi_d) = 
  \begin{bmatrix}
    \sin\theta_d\cos\varphi_d \\
    \sin\theta_d\sin\varphi_d \\
    -\cos\theta_d
  \end{bmatrix},
  \label{eq:deb_nd_sph}
\end{equation}
where $\theta_d \in [0,\pi/2]$ is the polar angle from nadir and 
$\varphi_d \in [0,2\pi)$ is the azimuth.  The tangent vectors to 
$\mathbb{S}^2$ at $\mathbf{n}_d$ are:
\begin{equation}
  \mathbf{u}_\theta = \frac{\partial\mathbf{n}_d}{\partial\theta_d}
  = \begin{bmatrix}
      \cos\theta_d\cos\varphi_d \\
      \cos\theta_d\sin\varphi_d \\
      \sin\theta_d
    \end{bmatrix},
  \qquad
  \mathbf{u}_\varphi = \frac{\partial\mathbf{n}_d}{\partial\varphi_d}
  = \begin{bmatrix}
      -\sin\theta_d\sin\varphi_d \\
      \sin\theta_d\cos\varphi_d \\
      0
    \end{bmatrix},
  \label{eq:deb_tangent}
\end{equation}
satisfying $\|\mathbf{u}_\theta\|=1$, 
$\|\mathbf{u}_\varphi\|=\sin\theta_d$, and 
$\mathbf{u}_\theta \cdot \mathbf{u}_\varphi = 0$.
The Jacobian from spherical to Cartesian is:
\begin{equation}
  \mathbf{J}_{\mathrm{sph}} 
  = [\,\mathbf{u}_\theta \;\; \mathbf{u}_\varphi\,]
  \in \mathbb{R}^{3\times 2}.
  \label{eq:deb_Jsph}
\end{equation}

% -----------------------------------------------------------------
\subsection*{3. Fisher Information Matrix}

The full parameter vector is 
$\boldsymbol{\xi} = [\theta_d,\,\varphi_d,\,\alpha]^{\mathsf{T}}$,
where $\alpha = \ln\eta$ is a log-reparameterization of the nuisance 
for numerical stability.  Under the Gaussian model~\eqref{eq:deb_obs}, the 
FIM follows from the Slepian--Bangs formula~\cite{Kay1993}:
\begin{equation}
  \mathcal{I}(\boldsymbol{\xi})
  = \frac{N}{\sigma^2}
    \sum_{i=1}^{K}
    \bigl[\nabla_{\boldsymbol{\xi}}\,\mu_i\bigr]\,
    \bigl[\nabla_{\boldsymbol{\xi}}\,\mu_i\bigr]^{\!\mathsf{T}},
  \label{eq:deb_FIM}
\end{equation}
where $\nabla_{\boldsymbol\xi} = 
[\partial/\partial\theta_d,\;\partial/\partial\varphi_d,\;
\partial/\partial\alpha]^{\mathsf{T}}$.

With $\mu_i = e^\alpha\,Q_i^m$ where $Q_i = \mathbf{n}_{t,i}\cdot\mathbf{n}_d$, 
the required partial derivatives are:
\begin{subequations}\label{eq:deb_grads}
\begin{align}
  \frac{\partial\mu_i}{\partial\theta_d}
  &= \eta\,m\,Q_i^{m-1}\,
     (\mathbf{n}_{t,i}\cdot\mathbf{u}_\theta),
  \label{eq:deb_grad_theta}\\
  \frac{\partial\mu_i}{\partial\varphi_d}
  &= \eta\,m\,Q_i^{m-1}\,
     (\mathbf{n}_{t,i}\cdot\mathbf{u}_\varphi),
  \label{eq:deb_grad_phi}\\
  \frac{\partial\mu_i}{\partial\alpha}
  &= \eta\,Q_i^m.
  \label{eq:deb_grad_alpha}
\end{align}
\end{subequations}
Factoring out $\eta$ and defining the scaled gradient
\begin{equation}
  \tilde{\mathbf{g}}_i = 
  \begin{bmatrix}
    m\,Q_i^{m-1}\,(\mathbf{n}_{t,i}\cdot\mathbf{u}_\theta) \\
    m\,Q_i^{m-1}\,(\mathbf{n}_{t,i}\cdot\mathbf{u}_\varphi) \\
    Q_i^m
  \end{bmatrix},
  \label{eq:deb_gtilde}
\end{equation}
the FIM becomes:
\begin{equation}
  \mathcal{I}(\boldsymbol{\xi})
  = \frac{N\eta^2}{\sigma^2}
    \sum_{i=1}^{K}
    \tilde{\mathbf{g}}_i\,\tilde{\mathbf{g}}_i^{\mathsf{T}}
  \;\in\;\mathbb{R}^{3\times3}.
  \label{eq:deb_FIM_compact}
\end{equation}

% -----------------------------------------------------------------
\subsection*{4. Direction Error Bound}

Since $\eta$ (equivalently $\alpha$) is a nuisance parameter, the CRLB 
for the angular parameters $(\theta_d,\varphi_d)$ is obtained from the 
inverse of the full FIM.  Partitioning:
\begin{equation}
  \mathcal{I} = 
  \begin{bmatrix}
    \mathbf{I}_{\mathrm{ang}} & \mathbf{i}_{\mathrm{cross}} \\
    \mathbf{i}_{\mathrm{cross}}^{\mathsf{T}} & I_{\alpha\alpha}
  \end{bmatrix},
\end{equation}
where $\mathbf{I}_{\mathrm{ang}} \in \mathbb{R}^{2\times2}$, 
$\mathbf{i}_{\mathrm{cross}} \in \mathbb{R}^{2\times1}$, and 
$I_{\alpha\alpha} \in \mathbb{R}$, the angular CRLB is:
\begin{equation}
  \operatorname{Cov}[\hat\theta_d,\hat\varphi_d]
  \;\succeq\;
  \mathbf{C}_{\mathrm{ang}}
  \;\triangleq\;
  \bigl[\mathcal{I}^{-1}\bigr]_{1:2,\,1:2}
  = \bigl(\mathbf{I}_{\mathrm{ang}} 
    - \mathbf{i}_{\mathrm{cross}}\,I_{\alpha\alpha}^{-1}\,
      \mathbf{i}_{\mathrm{cross}}^{\mathsf{T}}\bigr)^{-1}.
  \label{eq:deb_Cang}
\end{equation}

The angular covariance $\mathbf{C}_{\mathrm{ang}}$ is in the 
$(\theta_d,\varphi_d)$ coordinate system.  To express it in Cartesian 
direction-error space, we propagate through the Jacobian 
$\mathbf{J}_{\mathrm{sph}}$ from~\eqref{eq:deb_Jsph}:
\begin{equation}
  \operatorname{Cov}[\delta\mathbf{n}_d]
  \;\succeq\;
  \mathbf{J}_{\mathrm{sph}}\,\mathbf{C}_{\mathrm{ang}}\,
  \mathbf{J}_{\mathrm{sph}}^{\mathsf{T}}.
\end{equation}

The Direction Error Bound is the square root of the trace:
\begin{equation}
  \boxed{
  \mathrm{DEB}(\mathbf{n}_d)
  = \sqrt{\operatorname{tr}\!\bigl\{
    \mathbf{J}_{\mathrm{sph}}\,\mathbf{C}_{\mathrm{ang}}\,
    \mathbf{J}_{\mathrm{sph}}^{\mathsf{T}}
    \bigr\}}
  = \sqrt{
    \mathbf{C}_{\mathrm{ang}}(1,1) 
    + \sin^2\!\theta_d\;\mathbf{C}_{\mathrm{ang}}(2,2)
    }.
  }
  \label{eq:DEB}
\end{equation}
The DEB represents the minimum achievable RMS direction error 
$\sqrt{\mathbb{E}[\|\hat{\mathbf{n}}_d - \mathbf{n}_d\|^2]}$
for any unbiased direction estimator using the $K$ direction-finding 
measurements.  For small errors, $\mathrm{DEB} \approx$ angular RMSE 
in radians.

\emph{Remark.}
The DEB uses only the $K$ direction-finding measurements.  
The distance-recovery measurement ($K{+}1$) does not contribute because 
$\eta$ is treated as a nuisance parameter that is profiled out via the 
Schur complement.  This makes the DEB a fair benchmark for all 
direction estimators (GLS, WLS, NLS), which also use only $K$ measurements.

\bibliographystyle{IEEEtran}
\bibliography{References}
\end{document}
